\section{Conclusion}
\label{sec:conclusion}

Projection-blind delayed redistributions contain a sharp amount of recoverable
information.  Their zeroth delay moment vanishes exactly, but two separated
delay atoms allow the first moment to generate every transverse forcing
direction.  The source map has a sharp right inverse, and the stable Markov
inverse transfers this controllability to a dimension-uniform observation map
on transverse return covectors.  When the merged delay support contains only
one location, the same source is identically zero.  This
necessity--sufficiency dichotomy is the operator core of the paper.

The RFDE analysis turns that core into a nonlinear readout.  A
Dobrushin-controlled complete-history graph, selected one-sided traces, and a
finite-section estimate show that normalized fold-matching root increments
recover the return pairings with error \(C(\delta+|\zeta|)\), uniformly over
network size and the unit probe ball.  Generator-supported directions preserve
the chosen base-layer support, and special positive base layers admit a common
signed perturbation interval.  Heterogeneous fast curvature and heterogeneous slow
recovery sensing yield distinct return covectors; the second model class
realizes every nonzero transverse covector ray at fixed network size.  Sparse
reset--cycle families show that these effects persist on \(O(N)\)-edge
networks rather than only in dense mean-field examples.

The identification claim is deliberately narrower than full network
reconstruction.  The network skeleton, stationary weights, delay locations,
and scalar coefficients are known; the root responses identify the compressed
return covector.  Its observation map is uniformly conditioned as an
isomorphism onto its range when that range carries the inherited dual norm.
Coordinate recovery from a finite noisy probe frame needs a separate lower
frame bound, and absolute root-measurement noise
is amplified by \(\delta^{-3}|\zeta|^{-1}\).  Thus the present theorem gives
uniform mathematical identifiability of weighted pairings, not a noise-robust
experimental inversion theorem.

The root itself has an equally sharp scope boundary.  Localized admissible
changes of the artificial outer completion move the exact finite-\(\delta\)
root while leaving the physical RFDE and retained fold data fixed.  The current
finite-section axioms therefore do not define a preparation-independent
maximal canard.  Once each preparation's baseline is subtracted, however, the
leading response germ and its limiting covector are independent of the fixed
preparation within the stated comparison class.  The relative-canonical
object proved here is the hidden-return response functional with the physical
model, projection, matching data, and parameter normalization fixed; it is
not the finite-section representative.

Closing Dobrushin gaps, moving or proliferating delay atoms, infinite networks,
and order-one layer changes would require new uniform range estimates.
Likewise, an intrinsic maximal canard would require physical outer invariant
manifolds or asymptotic boundary conditions absent from the local selection.
The present result isolates the reusable part of the mechanism: exact
projection blindness, delay-moment controllability, transverse return, and
nonlinear complete-history readout.
